\par\bigskip
\hypertarget{native-384-ceres2}{}
\begin{center}
\pdfbookmark[2]{2. Die Störungen der Planeten}{ceres-002}[2.]\par
\medskip
\textls[120]{Die Störungen der Planeten.}
\end{center}
\[
\begin{gathered}
\left.\begin{array}{c}x,y,z\\x',y',z'\end{array}\right\}
\text{ Coordinaten des }
\left\{\begin{array}{l}\text{gestörten}\\\text{störenden}\end{array}\right\},\\[4pt]
\mu=\frac{\text{Masse des st. Pl.}}{\text{Masse der Sonne}}
\end{gathered}
\]
\[
\begin{gathered}
r=\BareRoot{(xx+yy+zz)},\qquad
r'r'=\BareRoot{(x'x'+y'y'+z'z')},\\
\mathfrak{R}=\text{Abst[and] d. beiden Planeten [von einander]}
\end{gathered}
\]
\noindent\(m,A\) mittlere [Länge] und [mittlerer] Abstand eines beliebigen Plan[eten von der Sonne]
\[
\begin{aligned}
R&=\frac{xx'+yy'+zz'}{r'^{3}}-\frac{1}{\mathfrak{R}}\\
 &=\frac{r}{r'r'}\cos(\text{Dist. pl. angul.})-\frac{1}{\mathfrak{R}}\mathpunct{\cdot}
\end{aligned}
\]
\SourceLeaf{385}
\par\medskip
\noindent\hspace*{\parindent}\textls[120]{Fundamentalgleichungen:}
\[
\begin{aligned}
0&=\frac{\mathrm{d}\mathrm{d}x}{\mathrm{d}m^{2}}+\frac{xA^{3}}{r^{3}}+\mu A^{3}\left(\frac{\mathrm{d}R}{\mathrm{d}x}\right)\\[4pt]
0&=\frac{\mathrm{d}\mathrm{d}y}{\mathrm{d}m^{2}}+\frac{yA^{3}}{r^{3}}+\mu A^{3}\left(\frac{\mathrm{d}R}{\mathrm{d}y}\right)\\[4pt]
0&=\frac{\mathrm{d}\mathrm{d}z}{\mathrm{d}m^{2}}+\frac{zA^{3}}{r^{3}}+\mu A^{3}\left(\frac{\mathrm{d}R}{\mathrm{d}z}\right)
\end{aligned}
\]
\[
\begin{aligned}
[1]\quad 0&=\frac{\mathrm{d}r^{2}+rr\mathrm{d}v^{2}}{\mathrm{d}m^{2}}-\frac{2A^{3}}{r}+\frac{A^{3}}{a}+2\mu A^{3}\int\partial R\\[6pt]
[2]\quad 0&=\frac{\mathrm{d}r^{2}+r\mathrm{d}\mathrm{d}r}{\mathrm{d}m^{2}}-\frac{\mathrm{d}r^{2}+rr\mathrm{d}v^{2}}{\mathrm{d}m^{2}}\\
&\qquad+\frac{A^{3}}{r}+\mu A^{3}r\left(\frac{\mathrm{d}R}{\mathrm{d}r}\right)\\[6pt]
[3]\quad 0&=\frac{r\mathrm{d}\mathrm{d}r+\mathrm{d}r^{2}}{\mathrm{d}m^{2}}-\frac{A^{3}}{r}+\frac{A^{3}}{a}\\
&\qquad+2\mu A^{3}\int\partial R+\mu A^{3}r\left(\frac{\mathrm{d}R}{\mathrm{d}r}\right)\mathpunct{\cdot}
\end{aligned}
\]
\par
\noindent\hspace*{\parindent}Wir setzen
\[
2\mu A^{3}\int\partial R+\mu A^{3}r\left(\frac{\mathrm{d}R}{\mathrm{d}r}\right)=Q;
\]
[der Ausdruck \(\partial R\) ist so zu verstehen, dass bei der Differentiation nur die Coordinaten des gestörten Planeten zu variiren sind.
\par
\noindent\hspace*{\parindent}Sind \(\delta r\) die Störungen von \(r\), so wird aus der Gleichung 3 mit Vernachlässigung der 2. Potenzen der störenden Masse]
\[
[4]\qquad 0=\frac{\mathrm{d}\mathrm{d}(r\delta r)}{\mathrm{d}m^{2}}+\frac{A^{3}}{r^{3}}r\delta r+Q.
\]
\par
\noindent\hspace*{\parindent}[Zur Einführung der wahren Länge \(v\) als unabhängige Veränderliche schreibt sich] die vollständige Gleichung also, \(\delta r=\rho\) gesetzt,
\[
\begin{aligned}
[5]\quad 0&=\frac{r\mathrm{d}\mathrm{d}\rho+2\mathrm{d}r\mathrm{d}\rho+\rho\mathrm{d}\mathrm{d}r}{\mathrm{d}m^{2}}\\
&\qquad-(r\mathrm{d}\rho+\rho\mathrm{d}r)\frac{\mathrm{d}\mathrm{d}m}{\mathrm{d}m^{3}}+\frac{A^{3}}{r^{3}}r\rho+Q.
\end{aligned}
\]
\par
\noindent\hspace*{\parindent}Nun ist, die Excentric[ität] des gest. Pl.\(=e\) und \(\BareRoot{(1-ee)}=\gamma\) gesetzt, [mit demselben Genauigkeitsgrade]
\[
rr\mathrm{d}v=\gamma A^{\frac32}a^{\frac12}\mathrm{d}m
\]
und, \(\mathrm{d}v\) als constant betrachtet,
\[
2r\mathrm{d}r\mathrm{d}v=\gamma A^{\frac32}a^{\frac12}\mathrm{d}\mathrm{d}m
\]
[also wird Gleichung 5]
\[
[6]\qquad 0=\frac{r\mathrm{d}\mathrm{d}\rho+\rho\mathrm{d}\mathrm{d}r-\frac{2\rho}{r}\mathrm{d}r^{2}}{\mathrm{d}m^{2}}+\frac{A^{3}}{r^{3}}r\rho+Q.
\]
\par
\noindent\hspace*{\parindent}Endlich ist [wieder mit dem gleichen Genauigkeitsgrade]
\SourceLeaf{386}
\[
\begin{gathered}
\frac{\mathrm{d}\mathrm{d}\dfrac1r}{\mathrm{d}v^{2}}+\frac1r=\frac1{a\gamma\gamma}\qquad\text{oder}\\[4pt]
\frac{\mathrm{d}v^{2}}{r}-\frac{\mathrm{d}\mathrm{d}r}{rr}+\frac{2\mathrm{d}r^{2}}{r^{3}}=\frac{\mathrm{d}v^{2}}{a\gamma\gamma}\mathpunct{\cdot}
\end{gathered}
\]
\par
\noindent\hspace*{\parindent}Also
\[
[7]\qquad 0=\frac{r\mathrm{d}\mathrm{d}\rho+r\rho\mathrm{d}v^{2}-\frac{rr\rho}{a\gamma\gamma}\mathrm{d}v^{2}}{\mathrm{d}m^{2}}+\frac{A^{3}}{r^{3}}r\rho+Q
\]
oder [da das dritte und vierte Glied rechter Hand sich fortheben]
\[
0=\frac{\mathrm{d}\mathrm{d}\rho}{\mathrm{d}v^{2}}+\rho+\frac{\mathrm{d}m^{2}}{\mathrm{d}v^{2}}\frac{Q}{r}
\]
oder
\[
[8]\qquad 0=\frac{\mathrm{d}\mathrm{d}\rho}{\mathrm{d}v^{2}}+\rho+\frac{r^{3}Q}{A^{3}a\gamma\gamma}\mathpunct{\cdot}
\]
\par
\noindent\hspace*{\parindent}Hieraus ergibt sich sogleich
\[
\rho=\cos v\int\frac{r^{3}Q}{A^{3}a\gamma\gamma}\sin v\,\mathrm{d}v-\sin v\int\frac{r^{3}Q}{A^{3}a\gamma\gamma}\cos v\,\mathrm{d}v
\]
oder [da \(rr\mathrm{d}v=\gamma aa\mathrm{d}(nt)\)]
\[
*\quad\rho=\left\{\begin{aligned}&\cos v\int rQ\sin v\,\mathrm{d}(nt)\\[-2pt]&-\sin v\int rQ\cos v\,\mathrm{d}(nt)\end{aligned}\right\}\frac{a}{A^{3}\BareRoot{(1-ee)}}\mathpunct{\cdot}
\]
\par\medskip
\begin{center}\rule{.27\linewidth}{.4pt}\end{center}
\par
\noindent\hspace*{\parindent}[Ferner folgt aus den Gleichungen 1 und 3]
\[
\begin{aligned}
0&=\frac{\frac32\mathrm{d}r^{2}-\frac12rr\mathrm{d}v^{2}+2r\mathrm{d}\mathrm{d}r}{\mathrm{d}m^{2}}-\frac{A^{3}}{r}+\frac32\frac{A^{3}}{a}\\
&\qquad+3\mu A^{3}\int\partial R+2\mu A^{3}r\left(\frac{\mathrm{d}R}{\mathrm{d}r}\right)
\end{aligned}
\]
[also nach Einführung von \(\rho\) und \(\delta v\), Forthebung der ungestörten Werthe, und Vernachlässigung der 2. Potenzen der Störungen]
\[
\begin{aligned}
0&=\frac{3\mathrm{d}r\mathrm{d}\rho-rr\mathrm{d}v\mathrm{d}(\delta v)-r\rho\mathrm{d}v^{2}+2\rho\mathrm{d}\mathrm{d}r+2r\mathrm{d}\mathrm{d}\rho}{\mathrm{d}m^{2}}\\
&\qquad+\frac{A^{3}\rho}{rr}+3\mu A^{3}\int\partial R+2\mu A^{3}r\left(\frac{\mathrm{d}R}{\mathrm{d}r}\right)\\[6pt]
[9]\quad&=\frac{\mathrm{d}(\rho\mathrm{d}r+2r\mathrm{d}\rho)-rr\mathrm{d}v\mathrm{d}(\delta v)-r\rho\mathrm{d}v^{2}+\rho\mathrm{d}\mathrm{d}r}{\mathrm{d}m^{2}}\\
&\qquad+\frac{A^{3}\rho}{rr}+3\mu A^{3}\int\partial R+2\mu A^{3}r\left(\frac{\mathrm{d}R}{\mathrm{d}r}\right)\mathpunct{\cdot}
\end{aligned}
\]
\par
\noindent\hspace*{\parindent}Nun ist nach obiger Gleichung [2 mit ausreichender Genauigkeit]
\[
0=\frac{-r\mathrm{d}v^{2}+\mathrm{d}\mathrm{d}r}{\mathrm{d}m^{2}}+\frac{A^{3}}{rr}\mathpunct{\text{\raisebox{.45ex}{,}}}
\]
also
\[
\begin{aligned}
[10]\quad 0&=\frac{\mathrm{d}(\rho\mathrm{d}r+2r\mathrm{d}\rho)-rr\mathrm{d}v\mathrm{d}(\delta v)}{\mathrm{d}m^{2}}\\
&\qquad+3\mu A^{3}\int\partial R+2\mu A^{3}r\left(\frac{\mathrm{d}R}{\mathrm{d}r}\right)\mathpunct{\cdot}
\end{aligned}
\]
\par
\noindent\hspace*{\parindent}Folglich
\[
\begin{aligned}
*\quad\delta v&=\frac{\rho\mathrm{d}r+2r\mathrm{d}\rho}{aa\BareRoot{(1-ee)}\mathord{.}\mathrm{d}(nt)}\\
&\quad+\frac{3\mu a\int\!\int\partial R\mathord{.}\mathrm{d}(nt)+2\mu a\int r\left(\frac{\mathrm{d}R}{\mathrm{d}r}\right)\mathrm{d}(nt)}{\BareRoot{(1-ee)}}\mathpunct{\cdot}
\end{aligned}
\]
\par\medskip
\begin{center}\rule{.27\linewidth}{.4pt}\end{center}
\SourceLeaf{387}
\par
\noindent\hspace*{\parindent}Nehmen wir endlich für die Fundamentalebene, worin die Coordinaten \(x,y\) liegen, die ursprüngliche Bahn des [gestörten] Planeten an und setzen \(z=r\theta\), so ist
\[
0=\frac{\mathrm{d}\mathrm{d}(r\theta)}{\mathrm{d}m^{2}}+\frac{A^{3}\theta}{rr}+\mu A^{3}\left(\frac{\mathrm{d}R}{\mathrm{d}z}\right)\mathpunct{\text{\raisebox{.45ex}{,}}}
\]
welche [Gleichung] sich eben so integriren lässt wie obige [4], daher
\[
*\quad\theta=\left\{\begin{aligned}&\cos v\int r\left(\frac{\mathrm{d}R}{\mathrm{d}z}\right)\sin v\mathord{.}\mathrm{d}(nt)\\[2pt]&-\sin v\int r\left(\frac{\mathrm{d}R}{\mathrm{d}z}\right)\cos v\mathord{.}\mathrm{d}(nt)\end{aligned}\right\}\frac{\mu a}{\BareRoot{(1-ee)}}\mathpunct{\cdot}
\]
\par
\noindent\hspace*{\parindent}[\(\theta\) bedeutet die Breite resp. den Sinus der Breite über der als fest angenommenen ungestörten Bahnebene.] Dadurch wird also
\[
\begin{gathered}
\text{die hel[iocentrische] Breite vermehrt um }\\[2pt]
\frac{\cos(\mathrm{incl.})}{\cos(\mathrm{lat.\ hel.})}\theta\\[6pt]
\text{die hel[iocentrische] Länge vermindert um }\\[2pt]
\frac{\theta\sin(\mathrm{incl.})\cos(\mathrm{arg.\ lat.})}{\cos(\mathrm{lat.\ hel.})^{2}}\mathpunct{\cdot}
\end{gathered}
\]
\par\medskip
\begin{center}\rule{.27\linewidth}{.4pt}\end{center}
